
\documentclass{article}
\usepackage{colortbl}
\usepackage{makecell}
\usepackage{multirow}
\usepackage{supertabular}

\begin{document}

\newcounter{utterance}

\centering \large Interaction Transcript for game `cladder', experiment `full\_v1.5\_default', episode 7019 with refineonly{-}Qwen3.5{-}27B.
\vspace{24pt}

{ \footnotesize  \setcounter{utterance}{1}
\setlength{\tabcolsep}{0pt}
\begin{supertabular}{c@{$\;$}|p{.15\linewidth}@{}p{.15\linewidth}p{.15\linewidth}p{.15\linewidth}p{.15\linewidth}p{.15\linewidth}}
    \# & \multicolumn{2}{c}{Player} && \multicolumn{2}{c}{Game Master} \\
    \hline

    \theutterance \stepcounter{utterance}  
    & & & \multicolumn{4}{p{0.6\linewidth}}{
        \cellcolor[rgb]{0.9,0.9,0.9}{
            \makecell[{{p{\linewidth}}}]{
                \texttt{\tiny{[P1$\langle$GM]}}
                \texttt{You are an expert in causal inference. The following question is not a typical commonsense query, but rather a meticulously designed question created by a professor specializing in causal inference, intended to assess the students' mastery of the course content.} \\
\texttt{Imagine a self{-}contained, hypothetical world with only the following conditions, and without any unmentioned factors or causal relationships: Hwax has a direct effect on gyzp and jyka. Jyka has a direct effect on lirg. Gyzp has a direct effect on lirg. Method 1: We look at how jyka correlates with lirg case by case according to gyzp. Method 2: We look directly at how jyka correlates with lirg in general. To understand how jyka affects lirg, is it more correct to use the Method 1 than Method 2?} \\
\texttt{Start your answer with 'Yes' or 'No', followed by additional reasoning or evidence to support your explanation.} \\
            }
        }
    }
     \\ \\

    \theutterance \stepcounter{utterance}  
    & & \multicolumn{4}{p{0.6\linewidth}}{
        \cellcolor[rgb]{0.95,0.95,0.95}{
            \makecell[{{p{\linewidth}}}]{
                \texttt{\tiny{[P1$\leftrightarrow$P]}}
                \texttt{{[}DRAFT{]}} \\
\texttt{No} \\
\\ 
\texttt{In this hypothetical world, the causal structure is defined as follows: \$Hwax \textbackslash{}to Gyzp\$, \$Hwax \textbackslash{}to Jyka\$, \$Jyka \textbackslash{}to Lirg\$, and \$Gyzp \textbackslash{}to Lirg\$. This structure creates a specific confounding scenario. \$Hwax\$ acts as a common cause (confounder) for both \$Jyka\$ and \$Gyzp\$. Since \$Gyzp\$ is a direct cause of \$Lirg\$, \$Gyzp\$ functions as a **mediator** in the path from \$Hwax\$ to \$Lirg\$, but more importantly for the relationship between \$Jyka\$ and \$Lirg\$, \$Gyzp\$ is a **collider** on the path \$Jyka \textbackslash{}leftarrow Hwax \textbackslash{}to Gyzp \textbackslash{}to Lirg\$? No, let's re{-}evaluate the graph structure carefully.} \\
\\ 
\texttt{The paths are:} \\
\texttt{1.  \$Jyka \textbackslash{}to Lirg\$ (The direct causal effect we want to estimate).} \\
\texttt{2.  \$Jyka \textbackslash{}leftarrow Hwax \textbackslash{}to Gyzp \textbackslash{}to Lirg\$ (A backdoor path).} \\
\\ 
\texttt{In this graph, \$Hwax\$ is a confounder for the relationship between \$Jyka\$ and \$Lirg\$ because it causes both \$Jyka\$ and \$Gyzp\$ (which in turn causes \$Lirg\$). To identify the causal effect of \$Jyka\$ on \$Lirg\$, we must block the backdoor path \$Jyka \textbackslash{}leftarrow Hwax \textbackslash{}to Gyzp \textbackslash{}to Lirg\$.} \\
\\ 
\texttt{**Method 2** (looking at the general correlation) leaves this backdoor path open. The observed correlation will be a mix of the true causal effect of \$Jyka\$ on \$Lirg\$ and the spurious association induced by \$Hwax\$ flowing through \$Gyzp\$. Thus, Method 2 yields a biased estimate.} \\
\\ 
\texttt{**Method 1** involves conditioning on \$Gyzp\$ (looking at the correlation case by case according to \$Gyzp\$). By conditioning on \$Gyzp\$, we block the backdoor path \$Jyka \textbackslash{}leftarrow Hwax \textbackslash{}to Gyzp \textbackslash{}to Lirg\$. Since \$Gyzp\$ is a descendant of the confounder \$Hwax\$ and a direct cause of the outcome \$Lirg\$, conditioning on it effectively blocks the flow of non{-}causal association from \$Hwax\$ to \$Lirg\$ that is correlated with \$Jyka\$.} \\
\\ 
\texttt{Wait, let's double{-}check the "collider" logic.} \\
\texttt{Path: \$Jyka \textbackslash{}leftarrow Hwax \textbackslash{}to Gyzp \textbackslash{}to Lirg\$.} \\
\texttt{Here, \$Hwax\$ is a fork (common cause).} \\
            }
        }
    }
    & \\ \\

    \theutterance \stepcounter{utterance}  
    & \multicolumn{4}{p{0.6\linewidth}}{
        \cellcolor[rgb]{0.561,0.737,0.561}{
            \makecell[{{p{\linewidth}}}]{
                \texttt{\tiny{[P1$\rangle$GM]}}
                \texttt{No} \\
\\ 
\texttt{In this hypothetical world, the causal structure is defined as: \$Hwax \textbackslash{}to Gyzp\$, \$Hwax \textbackslash{}to Jyka\$, \$Jyka \textbackslash{}to Lirg\$, and \$Gyzp \textbackslash{}to Lirg\$. To determine the causal effect of \$Jyka\$ on \$Lirg\$, we must identify and block any "backdoor paths" that create spurious associations between the two.} \\
\\ 
\texttt{There is one such backdoor path: \$Jyka \textbackslash{}leftarrow Hwax \textbackslash{}to Gyzp \textbackslash{}to Lirg\$.} \\
\texttt{1.  **Method 2** (looking at the general correlation) leaves this path open. Because \$Hwax\$ influences both \$Jyka\$ and \$Gyzp\$ (which influences \$Lirg\$), the observed correlation between \$Jyka\$ and \$Lirg\$ will be confounded by \$Hwax\$. This results in a biased estimate of the causal effect.} \\
\texttt{2.  **Method 1** involves conditioning on \$Gyzp\$ (stratifying by \$Gyzp\$). In the backdoor path \$Jyka \textbackslash{}leftarrow Hwax \textbackslash{}to Gyzp \textbackslash{}to Lirg\$, \$Gyzp\$ is a non{-}collider node (specifically, it is a mediator in the path from \$Hwax\$ to \$Lirg\$, but a non{-}collider in the path connecting \$Jyka\$ and \$Lirg\$). Conditioning on a non{-}collider blocks the flow of association along that path. By fixing \$Gyzp\$, we sever the link between \$Hwax\$ and \$Lirg\$ in this specific chain, effectively blocking the backdoor path.} \\
\\ 
\texttt{Since Method 1 successfully blocks the confounding path while Method 2 does not, Method 1 provides an unbiased estimate of the causal effect, making it the more correct approach.} \\
            }
        }
    }
    & & \\ \\

    \theutterance \stepcounter{utterance}  
    & & & \multicolumn{2}{p{0.3\linewidth}}{
        \cellcolor[rgb]{0.9,0.9,0.9}{
            \makecell[{{p{\linewidth}}}]{
                \texttt{\tiny{[GM$|$GM]}}
                \texttt{No} \\
            }
        }
    }
    & & \\ \\

    \theutterance \stepcounter{utterance}  
    & & & \multicolumn{2}{p{0.3\linewidth}}{
        \cellcolor[rgb]{0.9,0.9,0.9}{
            \makecell[{{p{\linewidth}}}]{
                \texttt{\tiny{[GM$|$GM]}}
                \texttt{yes} \\
            }
        }
    }
    & & \\ \\

    \theutterance \stepcounter{utterance}  
    & & & \multicolumn{2}{p{0.3\linewidth}}{
        \cellcolor[rgb]{0.9,0.9,0.9}{
            \makecell[{{p{\linewidth}}}]{
                \texttt{\tiny{[GM$|$GM]}}
                \texttt{game\_result = LOSE} \\
            }
        }
    }
    & & \\ \\

\end{supertabular}
}

\end{document}
